% Options for packages loaded elsewhere
% Options for packages loaded elsewhere
\PassOptionsToPackage{unicode}{hyperref}
\PassOptionsToPackage{hyphens}{url}
\PassOptionsToPackage{dvipsnames,svgnames,x11names}{xcolor}
%
\documentclass[
  spanish,
  a4paper,
  oneside]{scrbook}
\usepackage{xcolor}
\usepackage[top=25mm,bottom=25mm,left=30mm,right=20mm]{geometry}
\usepackage{amsmath,amssymb}
\setcounter{secnumdepth}{5}
\usepackage{iftex}
\ifPDFTeX
  \usepackage[T1]{fontenc}
  \usepackage[utf8]{inputenc}
  \usepackage{textcomp} % provide euro and other symbols
\else % if luatex or xetex
  \usepackage{unicode-math} % this also loads fontspec
  \defaultfontfeatures{Scale=MatchLowercase}
  \defaultfontfeatures[\rmfamily]{Ligatures=TeX,Scale=1}
\fi
\usepackage{lmodern}
\ifPDFTeX\else
  % xetex/luatex font selection
\fi
% Use upquote if available, for straight quotes in verbatim environments
\IfFileExists{upquote.sty}{\usepackage{upquote}}{}
\IfFileExists{microtype.sty}{% use microtype if available
  \usepackage[]{microtype}
  \UseMicrotypeSet[protrusion]{basicmath} % disable protrusion for tt fonts
}{}
\makeatletter
\@ifundefined{KOMAClassName}{% if non-KOMA class
  \IfFileExists{parskip.sty}{%
    \usepackage{parskip}
  }{% else
    \setlength{\parindent}{0pt}
    \setlength{\parskip}{6pt plus 2pt minus 1pt}}
}{% if KOMA class
  \KOMAoptions{parskip=half}}
\makeatother
% Make \paragraph and \subparagraph free-standing
\makeatletter
\ifx\paragraph\undefined\else
  \let\oldparagraph\paragraph
  \renewcommand{\paragraph}{
    \@ifstar
      \xxxParagraphStar
      \xxxParagraphNoStar
  }
  \newcommand{\xxxParagraphStar}[1]{\oldparagraph*{#1}\mbox{}}
  \newcommand{\xxxParagraphNoStar}[1]{\oldparagraph{#1}\mbox{}}
\fi
\ifx\subparagraph\undefined\else
  \let\oldsubparagraph\subparagraph
  \renewcommand{\subparagraph}{
    \@ifstar
      \xxxSubParagraphStar
      \xxxSubParagraphNoStar
  }
  \newcommand{\xxxSubParagraphStar}[1]{\oldsubparagraph*{#1}\mbox{}}
  \newcommand{\xxxSubParagraphNoStar}[1]{\oldsubparagraph{#1}\mbox{}}
\fi
\makeatother


\usepackage{longtable,booktabs,array}
\newcounter{none} % for unnumbered tables
\usepackage{calc} % for calculating minipage widths
% Correct order of tables after \paragraph or \subparagraph
\usepackage{etoolbox}
\makeatletter
\patchcmd\longtable{\par}{\if@noskipsec\mbox{}\fi\par}{}{}
\makeatother
% Allow footnotes in longtable head/foot
\IfFileExists{footnotehyper.sty}{\usepackage{footnotehyper}}{\usepackage{footnote}}
\makesavenoteenv{longtable}
\usepackage{graphicx}
\makeatletter
\newsavebox\pandoc@box
\newcommand*\pandocbounded[1]{% scales image to fit in text height/width
  \sbox\pandoc@box{#1}%
  \Gscale@div\@tempa{\textheight}{\dimexpr\ht\pandoc@box+\dp\pandoc@box\relax}%
  \Gscale@div\@tempb{\linewidth}{\wd\pandoc@box}%
  \ifdim\@tempb\p@<\@tempa\p@\let\@tempa\@tempb\fi% select the smaller of both
  \ifdim\@tempa\p@<\p@\scalebox{\@tempa}{\usebox\pandoc@box}%
  \else\usebox{\pandoc@box}%
  \fi%
}
% Set default figure placement to htbp
\def\fps@figure{htbp}
\makeatother


% definitions for citeproc citations
\NewDocumentCommand\citeproctext{}{}
\NewDocumentCommand\citeproc{mm}{%
  \begingroup\def\citeproctext{#2}\cite{#1}\endgroup}
\makeatletter
 % allow citations to break across lines
 \let\@cite@ofmt\@firstofone
 % avoid brackets around text for \cite:
 \def\@biblabel#1{}
 \def\@cite#1#2{{#1\if@tempswa , #2\fi}}
\makeatother
\newlength{\cslhangindent}
\setlength{\cslhangindent}{1.5em}
\newlength{\csllabelwidth}
\setlength{\csllabelwidth}{3em}
\newenvironment{CSLReferences}[2] % #1 hanging-indent, #2 entry-spacing
 {\begin{list}{}{%
  \setlength{\itemindent}{0pt}
  \setlength{\leftmargin}{0pt}
  \setlength{\parsep}{0pt}
  % turn on hanging indent if param 1 is 1
  \ifodd #1
   \setlength{\leftmargin}{\cslhangindent}
   \setlength{\itemindent}{-1\cslhangindent}
  \fi
  % set entry spacing
  \setlength{\itemsep}{#2\baselineskip}}}
 {\end{list}}
\usepackage{calc}
\newcommand{\CSLBlock}[1]{\hfill\break\parbox[t]{\linewidth}{\strut\ignorespaces#1\strut}}
\newcommand{\CSLLeftMargin}[1]{\parbox[t]{\csllabelwidth}{\strut#1\strut}}
\newcommand{\CSLRightInline}[1]{\parbox[t]{\linewidth - \csllabelwidth}{\strut#1\strut}}
\newcommand{\CSLIndent}[1]{\hspace{\cslhangindent}#1}

\ifLuaTeX
\usepackage[bidi=basic,shorthands=off]{babel}
\else
\usepackage[bidi=default,shorthands=off]{babel}
\fi
\ifLuaTeX
  \usepackage{selnolig} % disable illegal ligatures
\fi


\setlength{\emergencystretch}{3em} % prevent overfull lines

\providecommand{\tightlist}{%
  \setlength{\itemsep}{0pt}\setlength{\parskip}{0pt}}



 


\AtBeginDocument{\hypersetup{linkcolor=black}}
% ---------- Paquetes tipograficos y microajustes ----------
\usepackage{microtype}        % mejor interletrado y justificado
\usepackage{csquotes}         % comillas tipograficas
\usepackage{iftex}
\ifPDFTeX\else
  % Seleccion de fuentes solo cuando XeTeX/LuaTeX esta disponible.
  \IfFontExistsTF{TeX Gyre Termes}{
    \setmainfont{TeX Gyre Termes}
  }{
    \IfFontExistsTF{Latin Modern Roman}{\setmainfont{Latin Modern Roman}}{}
  }
  \IfFontExistsTF{TeX Gyre Heros}{
    \setsansfont{TeX Gyre Heros}
  }{
    \IfFontExistsTF{Latin Modern Sans}{\setsansfont{Latin Modern Sans}}{}
  }
  \IfFontExistsTF{TeX Gyre Cursor}{
    \setmonofont{TeX Gyre Cursor}
  }{
    \IfFontExistsTF{Latin Modern Mono}{\setmonofont{Latin Modern Mono}}{}
  }
\fi
\usepackage{enumitem}         % listas compactas
\setlist{itemsep=.2em, topsep=.2em}

% ---------- KOMA-Script: estilo de titulos y espaciado ----------
\KOMAoptions{
  headings=big,
  parskip=half,
  fontsize=12pt,
  appendixprefix=true
}
\setlength{\parindent}{1.5em}
\setlength{\parskip}{0.6em}
\usepackage{indentfirst}
\usepackage{setspace}         % Control del interlineado del documento
\setstretch{1.5}

% ---------- Encabezados y pies (scrlayer-scrpage) ----------
\usepackage[automark,headsepline]{scrlayer-scrpage}
\clearpairofpagestyles
\automark[chapter]{chapter}
\ihead{\pagemark}
\ohead{\itshape\headmark}
\setheadsepline{0.4pt}
\renewcommand*{\chaptermarkformat}{}
\renewcommand*{\chapterpagestyle}{scrheadings}
\pagestyle{scrheadings}
\makeatletter
\let\ps@plain\ps@scrheadings
\makeatother
\cfoot{}

% ---------- Hipervinculos mas sobrios ----------
\usepackage{hyperref}
\hypersetup{
  colorlinks=true,
  linkcolor=black,
  citecolor=[rgb]{0.05,0.2,0.5},
  urlcolor=[rgb]{0.05,0.2,0.5},
  pdfauthor={\@author},
  pdftitle={\@title}
}

% ---------- Leyendas de figuras/tablas ----------
\usepackage[labelfont=bf,textfont=it]{caption}
\captionsetup{
  skip=8pt
}

% ---------- Tabla de contenidos ----------
\KOMAoptions{toc=graduated}
\RedeclareSectionCommand[tocnumwidth=3em]{chapter}
\RedeclareSectionCommand[tocindent=3.25em,tocnumwidth=2.8em]{section}
\RedeclareSectionCommand[tocindent=6.5em,tocnumwidth=2.5em]{subsection}
\makeatletter
\newcommand*{\tocdotfill}{\leavevmode\leaders\hbox to .6em{\hss.\hss}\hfill}
\makeatother
\RedeclareSectionCommand[toclinefill=\tocdotfill]{chapter}
\RedeclareSectionCommand[toclinefill=\tocdotfill]{section}
\RedeclareSectionCommand[toclinefill=\tocdotfill]{subsection}
\renewcommand*{\contentsname}{Tabla de contenido}
\setkomafont{chapterentry}{\normalfont}
\setkomafont{chapterentrypagenumber}{\normalfont}

% ---------- Viudas/Huerfanas y cortes de pagina ----------
\clubpenalty=10000
\widowpenalty=10000
\displaywidowpenalty=10000

% ---------- Entorno abstract para scrbook ----------
\providecommand{\abstractname}{Resumen}
\makeatletter
\@ifundefined{abstract}{
  \newenvironment{abstract}{
    \cleardoublepage
    \thispagestyle{plain}
    \null\vfill
    \begin{center}
      {\bfseries\Large \abstractname\par}
    \end{center}\vspace{1em}
    \begingroup
  }{
    \par\endgroup
    \vfill\null
    \cleardoublepage
  }
}{}
\makeatother

% ---------- Soporte de subtitulo desde YAML ----------
\makeatletter
\providecommand{\subtitle}[1]{\gdef\@subtitle{#1}}
\providecommand{\@subtitle}{}
\makeatother

\makeatletter
\newcommand{\facso@iflist}[4]{%
  \begingroup
    \def\addvspace##1{}%
    \def\numberline##1{##1}%
    \def\facso@target{#2}%
    \newcount\facso@entrycount
    \facso@entrycount=0\relax
    \long\def\contentsline##1##2##3{%
      \def\facso@this{##1}%
      \ifx\facso@this\facso@target
        \advance\facso@entrycount\@ne
      \fi
    }%
    \InputIfFileExists{\jobname.#1}{}{}%
  \endgroup
  \ifnum\facso@entrycount>0\relax
    #3%
  \else
    #4%
  \fi}
\renewcommand{\listoffigures}{%
  \facso@iflist{lof}{figure}{%
    \cleardoublepage
    \phantomsection
    \addchap*{Listado de Figuras}%
    \thispagestyle{scrheadings}%
    \markboth{Listado de Figuras}{Listado de Figuras}%
    \addcontentsline{toc}{chapter}{Listado de Figuras}%
    \@starttoc{lof}%
    \cleardoublepage
    \automark[chapter]{chapter}%
  }{}}
\renewcommand{\listoftables}{%
  \facso@iflist{lot}{table}{%
    \cleardoublepage
    \phantomsection
    \addchap*{Listado de Tablas}%
    \thispagestyle{scrheadings}%
    \markboth{Listado de Tablas}{Listado de Tablas}%
    \addcontentsline{toc}{chapter}{Listado de Tablas}%
    \@starttoc{lot}%
    \cleardoublepage
    \automark[chapter]{chapter}%
  }{}}
\makeatother

% --- Desactivar portada y abstract automaticos de Pandoc (PDF) ---
\AtBeginDocument{\let\maketitle\relax}
\renewenvironment{abstract}{}{}
\providecommand{\appendixname}{}
\providecommand{\appendixtocname}{}
\providecommand{\appendixpagename}{}
\renewcommand*{\appendixname}{Anexo}
\renewcommand*{\appendixtocname}{Anexos}
\renewcommand*{\appendixpagename}{Anexos}
\makeatletter
\@ifpackageloaded{bookmark}{}{\usepackage{bookmark}}
\makeatother
\makeatletter
\@ifpackageloaded{caption}{}{\usepackage{caption}}
\AtBeginDocument{%
\ifdefined\contentsname
  \renewcommand*\contentsname{Tabla de contenidos}
\else
  \newcommand\contentsname{Tabla de contenidos}
\fi
\ifdefined\listfigurename
  \renewcommand*\listfigurename{Listado de Figuras}
\else
  \newcommand\listfigurename{Listado de Figuras}
\fi
\ifdefined\listtablename
  \renewcommand*\listtablename{Listado de Tablas}
\else
  \newcommand\listtablename{Listado de Tablas}
\fi
\ifdefined\figurename
  \renewcommand*\figurename{Lista de figuras}
\else
  \newcommand\figurename{Lista de figuras}
\fi
\ifdefined\tablename
  \renewcommand*\tablename{Lista de tablas}
\else
  \newcommand\tablename{Lista de tablas}
\fi
}
\@ifpackageloaded{float}{}{\usepackage{float}}
\floatstyle{ruled}
\@ifundefined{c@chapter}{\newfloat{codelisting}{h}{lop}}{\newfloat{codelisting}{h}{lop}[chapter]}
\floatname{codelisting}{Listado}
\newcommand*\listoflistings{\listof{codelisting}{Listado de Listados}}
\makeatother
\makeatletter
\makeatother
\makeatletter
\@ifpackageloaded{caption}{}{\usepackage{caption}}
\@ifpackageloaded{subcaption}{}{\usepackage{subcaption}}
\makeatother
\usepackage{bookmark}
\IfFileExists{xurl.sty}{\usepackage{xurl}}{} % add URL line breaks if available
\urlstyle{same}
\hypersetup{
  pdftitle={Hacinamiento y saneamiento en Chile{,} un análisis comparativo entre zonas rurales y urbanas},
  pdfauthor={Violeta Mansilla; Francisca Osorio; Martina Soto},
  pdflang={es},
  colorlinks=true,
  linkcolor={Maroon},
  filecolor={Maroon},
  citecolor={Blue},
  urlcolor={Blue},
  pdfcreator={LaTeX via pandoc}}


\title{Hacinamiento y saneamiento en Chile, un análisis comparativo
entre zonas rurales y urbanas}
\usepackage{etoolbox}
\makeatletter
\providecommand{\subtitle}[1]{% add subtitle to \maketitle
  \apptocmd{\@title}{\par {\large #1 \par}}{}{}
}
\makeatother
\subtitle{Trabajo final}
\author{Violeta Mansilla \and Francisca Osorio \and Martina Soto}
\date{30 de junio de 2026}
\begin{document}
\frontmatter
\maketitle

% --- title-pdf.tex personalizado para portada formal ---
\makeatletter
\providecommand{\subtitle}[1]{\gdef\@subtitle{#1}}
\providecommand{\@subtitle}{}
\providecommand{\frontmattercontext}{}
\providecommand{\advisorname}{}
\providecommand{\advisorlabel}{FACSO}
\providecommand{\frontmatterlocation}{}
\InputIfFileExists{includes/cover-config.tex}{}{}

\newcommand{\PrintTitle}{%
  {\sffamily\bfseries\fontsize{24pt}{28pt}\selectfont \@title\par}%
}
\newcommand{\PrintSubtitle}{%
  \begingroup
  \edef\temp{\detokenize{\@subtitle}}%
  \ifx\temp\empty\relax
    % sin subtitulo
  \else
    {\sffamily\bfseries\large \@subtitle\par}%
  \fi
  \endgroup
}
\newcommand{\PrintAuthor}{%
  \begingroup
  \def\and{\\}% \and separa autores con salto de linea (en vez del \and de \maketitle estandar, que requiere una tabla abierta)
  {\Large\bfseries \@author\par}%
  \endgroup
}
\newcommand{\PrintDate}{%
  {\small \@date\par}%
}
\makeatother

\begin{titlepage}
\thispagestyle{empty}
\begin{center}
\vspace*{10mm}

% Logo o imagen institucional

\includegraphics[width=0.35\textwidth]{assets/cover.png}\par
\vspace{12mm}

% Titulo y subtitulo
\PrintTitle
\vspace{6mm}
\PrintSubtitle

\vspace{22mm}
\begingroup
\edef\temp{\detokenize{\frontmattercontext}}%
\ifx\temp\empty\relax
  % sin contexto adicional
\else
  {\normalsize \frontmattercontext\par}
\fi
\endgroup

\vspace{18mm}
\PrintAuthor

\vspace{16mm}
\begingroup
\edef\temp{\detokenize{\advisorname}}%
\ifx\temp\empty\relax
  % sin tutor
\else
  \rule{0.45\textwidth}{0.4pt}\par
  {\small \advisorlabel\ \advisorname\par}
\fi
\endgroup

\vfill
\begingroup
\edef\temp{\detokenize{\frontmatterlocation}}%
\ifx\temp\empty\relax
  % sin ubicacion
\else
  {\small \frontmatterlocation\par}
\fi
\endgroup
\PrintDate

\end{center}
\end{titlepage}

% Preliminares (numeros romanos) y estilo simple
\frontmatter
\pagestyle{scrheadings}

\setcounter{tocdepth}{2} % (o 1/3 segun prefieras)
\tableofcontents
\listoffigures
\listoftables


\mainmatter
\bookmarksetup{startatroot}

\chapter*{Abstract}\label{abstract}
\addcontentsline{toc}{chapter}{Abstract}

\markboth{Abstract}{Abstract}

\small\textbf{Palabras clave:} Hacinamiento, Saneamiento, Zona rural,
Zona urbana, Desigualdad, Acceso. \normalsize

El presente trabajo analiza las diferencias en las condiciones de
hacinamiento y acceso a saneamiento básico entre hogares de zonas
urbanas y rurales en Chile, a partir de los datos de la Encuesta de
Caracterización Socioeconómica Nacional (CASEN) 2024. Adicionalemente,
se realiza un análisis exploratorio de dicha problematica a nivel de
región, lo que permite fortalecer las discusiones del proyecto.

\bookmarksetup{startatroot}

\chapter{Introducción}\label{introducciuxf3n}

La presente investigación nace de la motivación por estudiar las
disparidades socioeconómicas y territoriales entre las áreas urbanas y
rurales en Chile ya que estas continúan siendo un desafío estructural
para el desarrollo del país. Las brechas en la calidad de vida no solo
se manifiestan en los ingresos monetarios, sino de manera crítica en las
condiciones habitacionales y en el acceso a servicios básicos dentro del
hogar (\citeproc{ref-ducci2008}{Ducci, 2008}). En este sentido, este
estudio aborda esta problemática centrándose en dos dimensiones
fundamentales: el hacinamiento y el acceso a saneamiento básico.

Frente a esto, la investigación plantea la siguiente pregunta central:
¿Cuáles son las diferencias en las situaciones de hacinamiento y de
acceso a saneamiento básico de los hogares en las zonas rurales respecto
de las zonas urbanas en Chile?. Para responder, el objetivo general de
este estudio es analizar las diferencias en las condiciones de
hacinamiento y acceso a saneamiento básico entre hogares de zonas
urbanas y rurales en Chile, a partir de los datos de la Encuesta CASEN
2024. Y en este sentido, los objetivos específicos contemplan: describir
la proporción de hogares en situación de hacinamiento en zonas urbanas y
rurales de Chile; comparar el acceso a saneamiento adecuado entre
hogares urbanos y rurales; y evaluar si las diferencias observadas en
hacinamiento y saneamiento entre zonas urbanas y rurales son
estadísticamente significativas.

La contribución esperada de este proyecto abierto es aportar evidencia
empírica reciente que pueda contribuir a la discusión pública, así como
en la profundización de políticas habitacionales focalizadas
territorialmente. De este modo, la estructura del documento se organiza
de forma lógica para dar cuenta de este proceso: Se comienza con la
presente sección introductoria, continua con los antecedentes que
presenta el respaldo teorico y empirico de esta investigación,
posteriormente la presentación de las hipótesis, seguido de los métodos
y sus resultados, para finalmente finalizar con las conclusiones.

Para poder cumplir con las interrogantes y los objetivos, el proyecto se
organiza mediante una justificación teórica y empírica de la
problemática, un sección de hipótesis, una explicación de los métodos
utilizados, seguido de sus resultados y finalmente una discusión de los
hallazgos. Asimismo, se espera que esta investigación pueda contribuir a
caracterizar esta problemática y sirva de respaldo empírico para la toma
de decisiones en políticas públicas.

\bookmarksetup{startatroot}

\chapter{Antecedentes}\label{antecedentes}

Como se mencionó anteriormente, esta investigación surge del interés por
estudiar las desigualdades sociales latentes en el país. Para ello, se
analiza una de sus principales aristas: la crisis habitacional actual,
expresada concretamente en las situaciones de hacinamiento y el déficit
de saneamiento básico de los hogares.

Desde una perspectiva conceptual, la habitabilidad y el acceso a la
infraestructura sanitaria minima constituyen componentes esenciales de
la calidad residencial. En este sentido, D'alencon et~al.
(\citeproc{ref-dalencon2008}{2008}) definen la habitabilidad como las
cualidades de la vivienda, su entorno inmediato y el conjunto
habitacional, abarcando desde la estructura material hasta el uso y
confort percibido, factores que permiten de manera conjunta ``la
satisfacción de las necesidades de las personas que la habitan y su
desarrollo'' (p.~275). De este modo, en el contexto de las asimetrías de
Latinoamérica, la vivienda opera como ``una expresión espacial de las
desigualdades sociales, y a la vez generadora de otras relacionadas con
el entorno urbano en donde se ubica'' (\citeproc{ref-ortiz2021}{Ortiz \&
Soria, 2021}).

Según la Encuesta de Caracterización Socioeconómica Nacional (CASEN)
(\citeproc{ref-mdsf2024}{Ministerio de Desarrollo Social y Familia,
2024}), el hacinamiento en el área urbana alcanza un 4,2\% frente a un
3,3\% en el sector rural; por otro lado, el saneamiento deficitario
afecta drásticamente al 18,4\% de los hogares rurales si lo comparamos
con un 2,8\% en las zonas urbanas.

Estos fenómenos se encuentran directamente vinculados con el desarrollo
habitacional del cono sur. Según Trebilcock
(\citeproc{ref-trebilcock2024}{2024}), la crisis de la vivienda en
América Latina se ha extendido de forma sostenida, donde Chile ha
experimentado un aumento crítico en el déficit habitacional desde el año
2017. Esta situación recientemente mencionada se ve complejizada por el
diseño de instrumentos, debido a que, como menciona Vergara-Perucich
(\citeproc{ref-vergara2023}{2023}), la política habitacional chilena
basada en claves subsidiarias focalizadas ha generado una especie de
círculo vicioso porque, si bien se observan reducciones en aspectos como
el allegamiento y la calidad material, el déficit continúa incrementando
en factores más complejos, tales como el hacinamiento.

Esta expresión del déficit habitacional cambia según el territorio,
concentrando dinámicas más complejas en las áreas urbanas debido a
factores propios de la oferta y demanda. Desde la oferta, la deficiente
planificación urbana ha propiciado el aumento de núcleos urbanos
hiperdensos, como es el caso de Estación Central, donde existen zonas
con densidades bastante altas que superan los 15.000 habitantes por
hectárea en edificaciones de gran altura y baja calidad
(\citeproc{ref-serey2023}{Serey Morales, 2023}). Por el lado de la
demanda, según Razmilic (\citeproc{ref-razmilic2019}{2019}), los flujos
migratorios recientes se han localizado en su mayoría en zonas urbanas
dinámicas, en lugar de zonas rurales, buscando proximidad a los mercados
laborales. En este sentido, la población migrante opta mayoritariamente
por el mercado de arriendo, lo que intensifica el hacinamiento crítico.

En lo que respecta al saneamiento básico, la brecha entre lo urbano y lo
rural responde a lógicas institucionales y de mercado bastante
diferenciadas. Los entornos urbanos se encuentran marcados por altos
niveles de cobertura y eficiencia; por el otro lado, la carencia tiende
a confinarse en los sectores más empobrecidos y marginados
(\citeproc{ref-berrios2023}{Berríos, 2023}). Este fenómeno de exclusión
en las ciudades se ve potenciado por la gentrificación, donde las
mejoras en la infraestructura básica incrementan el valor de la
propiedad, desplazando a las familias de bajos ingresos y, en este
sentido, fomentando las disparidades socioeconómicas
(\citeproc{ref-berrios2023}{Berríos, 2023}).

En el análisis de Berríos (\citeproc{ref-berrios2023}{2023}) se refleja
la realidad de las zonas rurales y las dificultades que enfrentan
respecto a gestión y de cobertura técnica antes que puramente
socioeconomicas, por ejemplo, esta disparidad se intensifica por la
variabilidad macro-regional del país y las falencias operacionales de
los Programas de Agua Potable Rural (APR), expresadas en una desigual
distribución presupuestaria y una falta de inversión en zonas
semiconcentradas o dispersas.

El análisis de estas variables mediante la CASEN 2024 resulta
fundamental, ya que permite visibilizar la disparidad de fenómenos como
el hacinamiento y el saneamiento en zonas rurales y urbanas.

\bookmarksetup{startatroot}

\chapter{Hipótesis}\label{hipuxf3tesis}

La pregunta que este proyecto busca respondes es: ¿Cuáles son las
diferencias en las situaciones de hacinamiento y de acceso a saneamiento
básico de los hogares en las zonas rurales respecto de las zonas urbanas
en Chile?

En ese sentido, para poder responder a dicha interrogante se plantean
las siguientes hipótesis:

H1: La proporción de hogares en situación de hacinamiento (más de 2.5
personas por dormitorio) es significativamente mayor en zonas rurales
que en zonas urbanas. H2: El acceso a saneamiento adecuado (agua potable
y sistema de alcantarillado) es menor en zonas rurales que en las zonas
urbanas.

Dichas hipótesis fueron pre registradas en el repositorio de Open
Science Framework (OSF) del proyecto, con el fin de garantizar la
transparencia y reproducibilidad de la investigación. El repositorio se
encuentra disponible en: Fuentes et~al.
(\citeproc{ref-fuentes2026}{2026})

\bookmarksetup{startatroot}

\chapter{Métodos}\label{muxe9todos}

Este trabajo cuenta con un diseño comparativo y transversal basado en la
observación, para lo cual se utilizó la base de datos y libro de códigos
de la Encuesta Casen 2024. A través de un enfoque cuantitativo se
comparó la situación de hacinamiento y de acceso a saneamiento básico
entre hogares urbanos y rurales. La unidad de análisis es el hogar, la
cual se identificó mediante la variable ``folio''.

Las variables dependientes son ``ind\_hacina'' e ``ind\_san'', las
cuales son indicadores de hacinamiento y de saneamiento ya construidos y
validados por el Ministerio de Desarrollo Social y Familia, que se
recodificaron como variables binarias (``hacinado'' y
``saneamiento\_deficitario'') para facilitar la comparación entre zonas.
La variable independiente principal es la zona de residencia (``area'',
recodificada como ``zona'': urbano/rural). La variable ``region'' es
incluída como control y para análisis exploratorio, mientras que
``varstrat'', ``varunit'' y ``expr'' son variables técnicas del diseño
muestral, no analíticas. El identificador ``folio'' es utilizado para
asegurar que cada hogar se cuente una sola vez.

Puesto que se trabaja a nivel de hogar, primero se eliminó la
duplicación de registros y se conserva un único registro por folio.
Luego se seleccionaron las variables necesarias para el análisis (folio,
ind\_hacina, ind\_san, area, region, expr, varstrat, varunit),
descartando el resto de la base. Finalmente, se excluyen los hogares con
valor ``No sabe'' (-88) en ind\_hacina, ya que no aportan información
válida para clasificar el nivel de hacinamiento. No se aplican otros
filtros, dado que la pregunta de investigación busca una comparación
general a nivel nacional, sin acotar a un subgrupo poblacional
específico.

El proyecto se organizó mediante la estructura IPO (Input,
Procesamiento, Output). En ese sentido, la base de datos de Casen (2024)
se encuentra dentro de la carpeta Input; los archivos de procesamiento,
análisis y visualización de los datos, se encuentran en la carpeta de
Procesamiento; y finalmente, las imágenes de los gráficos creados están
en la carpeta Output.

Cabe destacar que debido a que la base de datos utilizada tiene un peso
de 1.50GB, superando considerablemente el límite de GB de Github, esta
no pudo ser subida al repositorio del proyecto, por lo que solo se
importó de manera local al template. Esto significó una complicación al
momento de subir el proyecto al repositorio de Github, por lo que se
utilizó la herramienta de Gitignore para así conservar la base de datos
solo de manera local.

\bookmarksetup{startatroot}

\chapter{Resultados}\label{resultados}

Los análisis descriptivos realizados permitieron caracterizar la
situación de hacinamiento y saneamiento tanto en las zonas rurales como
en las zonas urbanas; asimismo, los resultados permitieron contrastar
las hipótesis anteriormente registradas.

Respecto de la situación de hacinamiento (H1), para evaluar la situación
según la zona de residencia, se realizó una prueba de Chi cuadrado junto
una regresión logística, utilizando las variables de hacinamiento y zona
(es decir, rural o urbana). Los resultados de las pruebas arrojaron que
el porcentaje de hacinamiento en las zonas urbanas es de 4,2\%, mientras
que en la zona rural es de 3,3\%, todo esto con una diferencia
estadísticamente significativa (p = 0,0006). A su vez el análisis de
regresión logística muestra que los hogares rurales tienen un 22\% menos
probabilidad de estar hacinados que los hogares urbanos (OR = 0,78). Por
tanto, la hipótesis que planteaba que el hacinamiento sería mayor en
zonas rurales que en zonas urbanas, se rechaza.

A continuación se presenta el gráfico con los resultados:

\begin{figure}[H]

\caption{Figura 1. Porcentaje de hogares en situación de hacinamiento
según zona geográfica. CASEN 2024. Elaboración propia.}

{\centering \pandocbounded{\includegraphics[keepaspectratio]{05-resultados_files/figure-pdf/unnamed-chunk-1-1.pdf}}

}

\end{figure}%

La metodología utilizada para contrastar la hipótesis de saneamiento
(H2) fue exactamente la misma, en ese sentido, los resultados indicaron
que mientras en las zonas urbanas el saneamiento deficitario es de
2,8\%, en las zonas rurales este alcanza un 18,4\%, por lo que los
hogares rurales tiene una probabilidad casi ocho veces mayor de vivir en
estas condiciones. Dichos resultados permiten confirmar la hipótesis de
que el acceso a saneamiento adecuado es menor en zonas rurales que en
las zonas urbanas.

A continuación se presenta el gráfico con los resultados:

\begin{figure}[H]

\caption{Figura 2. Porcentaje de hogares con saneamiento deficitario
según zona geográfica. CASEN 2024. Elaboración propia.}

{\centering \pandocbounded{\includegraphics[keepaspectratio]{05-resultados_files/figure-pdf/unnamed-chunk-2-1.pdf}}

}

\end{figure}%

Sobre los alcances del presente proyecto, este permite caracterizar la
situación de hacinamiento y saneamiento a nivel nacional, distinguiendo
entre zonas rurales y urbanas mediante el uso de datos recientes y
atingentes a la problemática. En ese sentido, cabe precisar que las
hipótesis registradas consideran únicamente el desglose urbano-rural; no
obstante, con el fin de enriquecer el estudio, se incorporó un análisis
exploratorio a nivel regional. Este último apartado se plantea sólo como
un complemento para la discusión, ya que no formaba parte del registro
inicial, pero se considera que su inclusión permite evaluar si las
brechas urbano-rurales son homogéneas en todo el país o si varían según
la región, fortaleciendo así las discusiones del proyecto.

En ese sentido, el análisis exploratorio arrojó que las diferencias en
hacinamiento entre zonas urbanas y rurales son pequeñas y poco
consistentes a lo largo del país, lo que es coherente con los hallazgos
de la hipótesis 1. Cabe señalar que las regiones que arrojan una mayor
crisis de hacinamiento corresponden a la Región de Arica y Parinacota y
la de Tarapacá, lo que podría atribuirse al alto flujo migratorio de
aquellas regiones. Respecto del análisis de saneamiento por región, este
arrojó diferencias más latentes, debido a que regiones del sur como La
Araucanía o Los Ríos presentan un mayor déficit de saneamiento que
regiones mayormente urbanizadas, como la Región Metropolitana. Por lo
que, tal como se presenta en los gráficos anteriores, es posible
destacar que las diferencias entre el hacinamiento y el acceso a
saneamiento adecuado, no son homogéneas a lo largo de todas las
regiones.

A continuación se presentan los gráficos con los resultados del análisis
exploratorio por región:

\begin{figure}[H]

\caption{Figura 3. Porcentaje de hogares en situación de hacinamiento
por región y zona geográfica. CASEN 2024. Elaboración propia.}

{\centering \pandocbounded{\includegraphics[keepaspectratio]{05-resultados_files/figure-pdf/unnamed-chunk-3-1.pdf}}

}

\end{figure}%

\begin{figure}[H]

\caption{Figura 4. Porcentaje de hogares con saneamiento deficitario por
región y zona geográfica. CASEN 2024. Elaboración propia.}

{\centering \pandocbounded{\includegraphics[keepaspectratio]{05-resultados_files/figure-pdf/unnamed-chunk-4-1.pdf}}

}

\end{figure}%

Para finalizar este apartado de resultados, es posible señalar que
dentro de los límites del proyecto se encuentra que los índices de
hacinamiento y saneamiento poseen categorías de respuesta bastante
generales; en ese sentido, hacinamiento solo se mide en términos de
cantidad de personas que habitan una vivienda, mientras que saneamiento
solo considera si dichas condiciones son aceptables o deficitarias,
ignorando un análisis más profundo de la problemática. Pese a ello, los
índices son útiles para medir la problemática de manera general.

\bookmarksetup{startatroot}

\chapter{Conclusiones}\label{conclusiones}

El presente proyecto permitió analizar una de las aristas más relevantes
de la crisis habitacional atravesada en los últimos años; la
problemática del hacinamiento y el déficit de acceso a saneamiento
adecuado. El análisis realizado permitió caracterizar esta problemática
tanto en zonas rurales como urbanas, así como desagregar la información
según las regiones del país.

En relación a las hipótesis planteadas, los resultados mostraron
conclusiones diferentes. La hipótesis sobre hacinamiento (H1) fue
rechazada, evidenciando que este problema afecta en mayor medida a los
hogares urbanos (4,2\%) que a los rurales (3,3\%). La hipótesis sobre
saneamiento (H2) fue confirmada de manera contundente, mostrando que el
déficit de acceso a saneamiento adecuado es dramáticamente mayor en
zonas rurales (18,4\%) que en zonas urbanas (2,8\%), con una
probabilidad casi 8 veces mayor de tener saneamiento deficitario en
hogares rurales (OR = 7,75).

Por su parte, el análisis exploratorio por región permitió reforzar
estos hallazgos, evidenciando que las brechas urbano-rurales no son
homogéneas a lo largo del país. Factores como la urbanización, la
migración y la distribución histórica de infraestructura sanitaria
podrían estar relacionados con los patrones observados, aunque el diseño
observacional del estudio no permite establecer relaciones causales.

En ese sentido, los resultados son consistentes, por lo que permiten
responder a la pregunta inicial, identificando una crisis habitacional
estructural que afecta de manera diferenciada a los territorios urbanos
y rurales del país. A su vez, detallan las diferencias en las
situaciones de hacinamiento y de acceso a saneamiento básico de los
hogares en las zonas rurales respecto de las zonas urbanas en Chile.

Esto sugiere la necesidad de políticas públicas focalizadas
territorialmente, que aborden tanto el hacinamiento en zonas
metropolitanas como el déficit de saneamiento en zonas rurales,
particularmente en las regiones del sur de Chile donde las brechas son
más pronunciadas.

\bookmarksetup{startatroot}

\chapter{Referencias}\label{referencias}

\protect\phantomsection\label{refs}
\begin{CSLReferences}{1}{0}
\bibitem[\citeproctext]{ref-berrios2023}
Berríos, J. (2023). \emph{Un acercamiento a las disparidades
urbano-rural en acceso al agua potable y saneamiento en {Chile}}
{[}Mathesis, Pontificia Universidad Católica de Chile{]}.
\url{https://repositorio.uc.cl/handle/11534/75436}

\bibitem[\citeproctext]{ref-dalencon2008}
D'alencon, R., Justiniano, C., Márquez, F., \& Valderrama, C. (2008).
Parámetros y estándares de habitabilidad: calidad en la vivienda, el
entorno inmediato y el conjunto habitacional. \emph{Camino al
Bicentenario: Propuestas para Chile}, 271-304.
\url{https://www.researchgate.net/publication/237832171}

\bibitem[\citeproctext]{ref-ducci2008}
Ducci, M. (2008). La política habitacional como instrumento de
desintegración social: Efectos de una política de vivienda exitosa. En
\emph{Retos para la integración social de los pobres en {América
Latina}} (pp. 293-310). {CLACSO}.

\bibitem[\citeproctext]{ref-fuentes2026}
Fuentes, V. M., Osorio, F. A., \& Ortega, M. S. (2026).
\emph{Hacinamiento y saneamiento en {Chile}, un análisis comparativo
entre zonas rurales y urbanas}.
\url{https://doi.org/10.17605/OSF.IO/96A3W}

\bibitem[\citeproctext]{ref-mdsf2024}
Ministerio de Desarrollo Social y Familia. (2024). \emph{Encuesta de
Caracterización Socioeconómica Nacional {CASEN} 2024}.
\url{https://observatorio.ministeriodesarrollosocial.gob.cl/encuesta-casen-2024}

\bibitem[\citeproctext]{ref-ortiz2021}
Ortiz, F. E., \& Soria, K. Y. (2021). El derecho a la vivienda: Apuntes
para su reivindicación desde la habitabilidad. \emph{Quid 16. Revista
del Área de Estudios Urbanos}, \emph{15}, 218-233.
\url{https://www.redalyc.org/journal/5596/559672582012/html/}

\bibitem[\citeproctext]{ref-razmilic2019}
Razmilic, S. (2019). Inmigración, vivienda y territorio. En
\emph{Inmigración en {Chile}: Una mirada multidimensional} (pp.
101-147).
\url{https://static.cepchile.cl/uploads/cepchile/2022/09/libro_inmigracion_vivienda.pdf}

\bibitem[\citeproctext]{ref-serey2023}
Serey Morales, C. (2023). \emph{El rol del {Estado} en el aseguramiento
del derecho a una vivienda adecuada en {Chile}} {[}Mathesis, Universidad
de Chile{]}. \url{https://repositorio.uchile.cl/handle/2250/195811}

\bibitem[\citeproctext]{ref-trebilcock2024}
Trebilcock, M. (2024). \emph{Persistencia y futuro de los asentamientos
informales en {Chile}} (Serie Policy Papers CiPP sobre Políticas
Públicas). Centro Interdisciplinar de Políticas Públicas ({CiPP}),
Universidad Alberto Hurtado.
\url{https://cippuahurtado.cl/wp-content/uploads/2025/01/Policy-asentamientos_con_numeracion_final-vfinal.pdf}

\bibitem[\citeproctext]{ref-vergara2023}
Vergara-Perucich, J. (2023). Contribución a la medición de la eficiencia
del sistema habitacional subsidiario en {Chile}. \emph{{ACE}:
Architecture, City and Environment}, \emph{18}(52).
\url{https://doi.org/10.5821/ace.18.52.11682}

\end{CSLReferences}

\cleardoublepage
\phantomsection
\addcontentsline{toc}{part}{Apéndices}
\appendix

\chapter{Instrumentos de
levantamiento}\label{instrumentos-de-levantamiento}

Incluye cuestionarios, guias de entrevista, etc.

\chapter{Tablas adicionales}\label{tablas-adicionales}

Tablas extendidas, tests de robustez, etc.


\backmatter

% --- Cuerpo del libro (capítulos) en arábigos ---
\mainmatter
\pagestyle{scrheadings}

% --- (Opcional) Estilo del índice general ---
% \addtocontents{toc}{\protect\thispagestyle{plain}}


\end{document}
